\documentclass[11pt]{letter}

\usepackage{geometry}
\geometry{margin=1in}
\usepackage{setspace}
\usepackage{hyperref}

\signature{Decory Edwards}
\address{Decory Edwards \\ Johns Hopkins University \\ Department of Economics}

\begin{document}

\begin{letter}{Hiring Committee \\ Google}

\opening{Dear Hiring Committee,}

I am writing to apply for the Data Scientist position with Google. I am currently a PhD candidate in Economics at Johns Hopkins University. My training combines mathematical reasoning, statistical modeling, and data driven empirical analysis. I am motivated by problems that require careful formulation, disciplined experimentation, and validation using large scale data. Google’s emphasis on applying scientific methods to product development and improvement aligns closely with how I approach research and analysis.

My research agenda is at the intersection of wealth inequality and macroeconomics. In my job market paper, I build and estimate a structural heterogeneous agent model with returns that differ across households. The core contribution is to show how heterogeneity in returns and income risk jointly shape the wealth distribution and consumption behavior. Relative to closely related work, my framework performs better at matching the lower and middle portions of the wealth distribution. This difference matters quantitatively. Models that focus primarily on returns heterogeneity can fit the top of the wealth distribution closely but tend to generate too much wealth among the bottom half of households. In my framework, income uncertainty generates a precautionary saving motive that produces genuinely low wealth households. This leads to a realistic distribution of marginal propensities to consume and aggregate consumption responses that align with empirical estimates.

A key feature of my research is the use of causal inference and structured empirical methods to answer applied questions. I design empirical strategies, work with large and complex datasets, and implement statistical models to isolate meaningful variation. I regularly use regression based methods, panel data techniques, and hypothesis testing to evaluate models and interventions. This work requires attention to assumptions, data quality, and interpretability, skills that translate directly to experimentation and product analytics.

I am particularly interested in Google because of its scale and its commitment to rigorous data driven decision making. The opportunity to collaborate with engineers, analysts, and product teams to translate business questions into tractable analysis is especially appealing. I value environments where statistical reasoning informs real product decisions and where results are evaluated through measurable outcomes.

My economics training has provided a strong foundation in probability, optimization, and statistical inference. I am comfortable working with Python, R, SQL, and related tools to extract, clean, and analyze data. I also value clear communication and collaboration and routinely present analytical results to both technical and non technical audiences. These skills are integral to my research process and would be central to my contributions as a data scientist.

I am able to begin a full time role in 2026 and am enthusiastic about working in a Google office environment. I value in person collaboration while also appreciating focused analytical work. I see this role as an opportunity to apply my academic training in a setting where analysis directly improves products used by millions of people.

I would welcome the opportunity to contribute my skills and perspective to Google. Thank you for your time and consideration.

\closing{Sincerely,}

\end{letter}
\end{document}
